\documentclass[UTF8]{ctexart}
\usepackage{xeCJK}
\usepackage{geometry}
\geometry{left=2cm,right=2cm,top=2.5cm,bottom=2.5cm}
\setlength{\headheight}{12.7pt}

\usepackage{titlesec}
\titleformat{\section}
  {\normalfont\large\bfseries}
  {\thesection}
  {1em}
  {}
\titlespacing*{\section}
  {0pt}
  {3.5ex plus 1ex minus .2ex}
  {2.3ex plus .2ex}

\usepackage{amsmath}
\usepackage{amssymb}
\usepackage{amsthm}
\usepackage{enumitem}
\setlist[enumerate]{itemsep=0pt,topsep=0pt,parsep=0pt,leftmargin=0pt,itemindent=2.5em}
\setlist[itemize]{itemsep=0pt,topsep=0pt,parsep=0pt,leftmargin=0pt,itemindent=2.5em}
\usepackage{fancyhdr}
\pagestyle{fancy}
\fancyhf{}
\usepackage{graphicx}
\usepackage{hyperref}
\usepackage{indentfirst}
\usepackage{lmodern}


\begin{document}
\setlength{\abovedisplayskip}{1pt}
\setlength{\belowdisplayskip}{1pt}

\section{米田嵌入}

这里给定宇宙$\mathcal{U}$.

\vspace{10pt}

\hypertarget{sec:定义1.1}{}
\noindent\textbf{定义1.1 (预层范畴)}

给定\href{../01 范畴、函子与自然变换/01 范畴#nameddest=sec:定义1.2}{局部小范畴}
$\mathbf{C}$, 称从$\mathbf{C}^{\mathrm{op}}$到$\mathbf{Set}$的函子
为$\mathbf{C}$上的\textbf{预层}, 定义
\[
    \mathbf{PSh}(\mathbf{C})\overset{\mathrm{def}}{=}
    [\mathbf{C}^{\mathrm{op}},\mathbf{Set}],
\]
称之为$\mathbf{C}$上的\textbf{预层范畴}.

\vspace{10pt}

\hypertarget{sec:定义1.2}{}
\noindent\textbf{定义1.2 (表出预层)}

给定局部小范畴$\mathbf{C}$和其中的对象$X$,
\begin{enumerate}
    \item 对$\mathbf{C}^{\mathrm{op}}$中任意的对象$A$,
    定义$h_X(A)=\mathbf{C}(A,X)$,
    \item 对$\mathbf{C}^{\mathrm{op}}$中任意的态射$u:A\to B$,
    即$\mathbf{C}$中的态射$u:B\to A$, 定义映射
    \[
        h_X(u):\mathbf{C}(A,X)\to\mathbf{C}(B,X),\quad
        \alpha\mapsto\alpha\circ u,
    \]
\end{enumerate}
则$h_X$是$\mathbf{C}$上的预层, 称之为\textbf{由}$X$\textbf{表出的预层}.

\noindent{\kaishu 验证.}

\begin{enumerate}
    \item 对$\mathbf{C}^{\mathrm{op}}$中任意的对象$A$, 有
    \[\left\{\begin{aligned}
        &h_X(I_A):&&\alpha\mapsto\alpha\circ I_A=\alpha,\\[-3pt]
        &I_{h_X(A)}=\mathrm{id}_{\mathbf{C}(A,X)}:&&\alpha\mapsto\alpha,
    \end{aligned}\right.\]
    从而$h_X(I_A)=I_{h_X(A)}$.
    \item 对$\mathbf{C}^{\mathrm{op}}$中任意的态射$u,v$, 在有定义的情形下有
    \[\left\{\begin{aligned}
        &h_X(u\circ^{\mathrm{op}}v):&&\alpha\mapsto\alpha\circ(v\circ u),\\[-3pt]
        &h_X(u)\circ h_X(v):&&\alpha\mapsto(\alpha\circ v)\circ u,
    \end{aligned}\right.\]
    从而$h_X(u\circ^\mathrm{op}v)=h_X(u)\circ h_X(v)$. \hfill $\qedsymbol$
\end{enumerate}

\vspace{50pt}

\begin{center}\begin{tikzpicture}[
    box1/.style={draw, thick, minimum width=2cm, minimum height=3cm},
    box2/.style={draw, thick, minimum width=6cm, minimum height=3cm},
    arrow/.style={-Stealth, thick}]
    \node[box1] (Copbox) at (-7,0) {};
    \node[anchor=north west, xshift=5, yshift=-5] at (Copbox.north west)
    {$\mathbf{C}^{\mathrm{op}}$};
    \node (A) at (-7.7, 0) {$A$};
    \node (B) at (-6.3, 0) {$B$};
    \draw[arrow] (A) -- node[above] {$u$} (B);
    \node[box2] (Setbox) at (0,0) {};
    \node[anchor=north west, xshift=5, yshift=-5] at (Setbox.north west)
    {$\mathbf{Set}$};
    \node (CAX) at (-2, 0) {$\mathbf{C}(A,X)$};
    \node (CBX) at (2, 0) {$\mathbf{C}(B,X)$};
    \draw[arrow] (CAX) -- node[above] {右复合$u$} (CBX);
    \draw[arrow, very thick] ([xshift=10]Copbox.east)
    --node[above] {$h_X$} ([xshift=-10]Setbox.west);
    \node (A_) at (4, 0) {$A$};
    \node (X_) at (5.5, 1) {$X$};
    \node (B_) at (7, 0) {$B$};
    \draw[arrow] (A_) -- node[left, yshift=2] {$\alpha$} (X_);
    \draw[arrow] (B_) -- node[below] {$u$} (A_);
    \draw[arrow] (B_) -- node[right, yshift=2] {$\alpha\circ u$} (X_);
\end{tikzpicture}\end{center}





\newpage

\hypertarget{sec:定义1.3}{}
\noindent\textbf{定义1.3 (米田嵌入)}

给定局部小范畴$\mathbf{C}$, 对$\mathbf{C}$中任意的对象$X$可以由它表出预层$h_X$.

然后对$\mathbf{C}$中任意的态射$f:X\to Y$定义
\[
    h_f:\mathrm{Ob}(\mathbf{C}^{\mathrm{op}})\to\mathrm{Mor}(\mathbf{Set}),
\]
其中对$\mathbf{C}^{\mathrm{op}}$中任意的对象$A$,
\[
    h_f(A):\mathbf{C}(A,X)\to\mathbf{C}(A,Y),\quad
    \alpha\mapsto f\circ\alpha,
\]
则$h_f:h_X\Rightarrow h_Y$.

进一步, $h$是从$\mathbf{C}$到$\mathbf{PSh}(\mathbf{C})$的函子,
称之为\textbf{米田嵌入}.

\noindent{\kaishu 验证.}

首先验证$h_f:h_X\Rightarrow h_Y$. 对$\mathbf{C}^{\mathrm{op}}$中任意的态射
$u:A\to B$, 在$\mathbf{Set}$中有
\[\left\{\begin{aligned}
    &h_Y(u)\circ h_f(A):&&\alpha\mapsto(f\circ\alpha)\circ u,\\[-3pt]
    &h_f(B)\circ h_X(u):&&\alpha\mapsto f\circ(\alpha\circ u),
\end{aligned}\right.\]
从而$h_Y(u)\circ h_f(A)=h_f(B)\circ h_X(u)$, 即$h_f$是自然变换.

然后验证$h:\mathbf{C}\to\mathbf{PSh}(\mathbf{C})$.
\begin{enumerate}
    \item 对$\mathbf{C}$中任意的对象$X$和$A$, 有
    \[\left\{\begin{aligned}
        &h_{I(X)}(A):&&\alpha\mapsto I(X)\circ\alpha=\alpha,\\[-3pt]
        &I_{h_X}(A)=\mathrm{id}_{h_X(A)}:&&\alpha\mapsto\alpha,
    \end{aligned}\right.\]
    从而$h_{I(X)}=I_{h_X}$.
    \item 对$\mathbf{C}$中任意的态射$f,g$,
    在有定义的情形下对$\mathbf{C}$中任意的对象$A$有
    \[\left\{\begin{aligned}
        &h_{g\circ f}(A):&&\alpha\mapsto(g\circ f)\alpha,\\[-3pt]
        &(h_g\circ h_f)(A):&&\alpha\mapsto g\circ(f\circ\alpha),
    \end{aligned}\right.\]
    从而$h_{g\circ f}=h_g\circ h_f$. \hfill $\qedsymbol$
\end{enumerate}

\vspace{10pt}

\begin{center}\begin{tikzpicture}[
    box1/.style={draw, thick, minimum width=2cm, minimum height=5cm},
    box2/.style={draw, thick, minimum width=6cm, minimum height=5cm},
    arrow/.style={-Stealth, thick}]
    \node[box1] (Copbox) at (-7,0) {};
    \node[anchor=north west, xshift=5, yshift=-5] at (Copbox.north west)
    {$\mathbf{C}^{\mathrm{op}}$};
    \node (A) at (-7.7, 0) {$A$};
    \node (B) at (-6.3, 0) {$B$};
    \draw[arrow] (A) -- node[above] {$u$} (B);
    \node[box2] (Setbox) at (0,0) {};
    \node[anchor=north west, xshift=5, yshift=-5] at (Setbox.north west)
    {$\mathbf{Set}$};
    \node (CAX) at (-2, 1) {$\mathbf{C}(A,X)$};
    \node (CBX) at (2, 1) {$\mathbf{C}(B,X)$};
    \draw[arrow] (CAX) -- node[above] {右复合$u$} (CBX);
    \draw[arrow, very thick] ([xshift=10, yshift=1cm]Copbox.east)
    --node[above, pos=0.3] {$h_X$} ([xshift=-10, yshift=1cm]Setbox.west);
    \node (CAY) at (-2, -1) {$\mathbf{C}(A,Y)$};
    \node (CBY) at (2, -1) {$\mathbf{C}(B,Y)$};
    \draw[arrow] (CAY) -- node[above] {右复合$u$} (CBY);
    \draw[arrow, very thick] ([xshift=10, yshift=-1cm]Copbox.east)
    --node[above, pos=0.3] {$h_Y$} ([xshift=-10, yshift=-1cm]Setbox.west);
    \draw[arrow] (CAX) -- node[right] {左复合$f$} (CAY);
    \draw[arrow] (CBX) -- node[left] {左复合$f$} (CBY);
    \draw[double, double distance=2pt, thick, -Latex] (-4.5, 0.7)
    --node[right] {$h_f$} (-4.5, -0.7);
    \node (A_) at (4, 0) {$A$};
    \node (X_) at (5.5, 1) {$X$};
    \node (B_) at (7, 0) {$B$};
    \node (Y_) at (5.5, -1) {$Y$};
    \draw[arrow] (A_) -- node[left, yshift=2] {$\alpha$} (X_);
    \draw[arrow] (B_) -- node[below, pos=0.7] {$u$} (A_);
    \draw[arrow] (B_) -- node[right, yshift=2] {$\alpha\circ u$} (X_);
    \draw[arrow] (A_) -- node[left, yshift=-4] {$f\circ\alpha$} (Y_);
    \draw[arrow] (X_) -- node[right, pos=0.3] {$f$} (Y_);
    \draw[arrow] (B_) -- node[right, yshift=-4] {$f\circ\alpha\circ u$} (Y_);
\end{tikzpicture}\end{center}

\end{document}